% Circuito RLC em série, com fonte alternada
% Feito para o ÉguaTeX (eguatex.com.br). Uso livre (CC0 1.0).
%
% Compile e pronto: sai só a figura, em PDF.
% Para pôr no seu trabalho, use o PDF com \includegraphics ou copie
% o bloco circuitikz para dentro de uma figure e ponha no preâmbulo
% \usepackage[american]{circuitikz}.
% Para os símbolos europeus (resistor em retângulo), troque american por european.
\documentclass[border=4pt]{standalone}
\usepackage[T1]{fontenc}
\usepackage{lmodern}
\usepackage[american]{circuitikz}

\begin{document}
\begin{circuitikz}[line cap=round, line join=round]

% Mude estes textos e tamanhos.
\def\fonte{$v(t)$}         % fonte alternada
\def\resistor{$R$}         % resistor
\def\indutor{$L$}          % indutor
\def\capacitor{$C$}        % capacitor
\def\largura{6}            % largura do circuito, em cm
\def\altura{3.5}           % altura do circuito, em cm

\draw (0,0) to[sV, l=\fonte] (0,\altura)
            to[R, l=\resistor, i>^=$i(t)$] ({\largura/2},\altura)
            to[L, l=\indutor] (\largura,\altura)
            to[C, l=\capacitor] (\largura,0)
            -- (0,0);
\end{circuitikz}
\end{document}
